% ECONOMICS_PROBLEM: {"id": "C61-93", "title": "Discounted continuous-time turnpikes and the discount-dependent reference", "jel_code": "C61", "source_id": "OP-0572"}
% FIRST_POSTED: 2026-10-10
% ATTEMPT_STATUS: solved
% ENTRY_KIND: research_attempt
% AUTHOR: Ji Ho Bae
% AI_ASSISTANCE: GPT Codex; internal symbolic and scope audits; no Lean formalization; no external peer review
% SOLUTION_SCOPE: Sufficient structural answer to the formal agenda
% NOVELTY_CLAIM: none
% INTERNAL_REVIEW_RECORD: ../CHILD_REVIEW.md
\documentclass[11pt,a4paper]{article}
\usepackage[margin=25mm]{geometry}
\usepackage{amsmath,amssymb,amsthm,mathtools,hyperref}
\newtheorem{theorem}{Theorem}
\newtheorem{lemma}[theorem]{Lemma}
\newtheorem{proposition}[theorem]{Proposition}
\newtheorem{corollary}[theorem]{Corollary}
\newtheorem{remark}[theorem]{Remark}
\newcommand{\R}{\mathbb R}
\newcommand{\Leb}{\operatorname{Leb}}
\newcommand{\Lip}{\operatorname{Lip}}
\newcommand{\diam}{\operatorname{diam}}
\newcommand{\argmin}{\operatorname*{arg\,min}}
\newcommand{\sym}{\operatorname{sym}}
\newcommand{\logplus}{\log_{+}}
\title{A verifiable nonlinear sufficient-condition theorem for discounted continuous-time turnpikes}
\author{Ji Ho Bae}
\date{10 October 2026}
\begin{document}
\maketitle

\begin{abstract}
We give a sufficient structural answer to the formal agenda of C61-93.
For arbitrary finite-dimensional systems with a nonlinear contractive drift,
compact convex control constraints, a strongly convex control cost, and a
possibly nonconvex state cost, explicit bounds on the problem data construct
a unique discount-dependent stationary reference. Every finite-horizon
optimizer satisfies an unweighted exponential turnpike estimate, uniformly
over the stated initial-state and terminal-cost family. We give the
constants and rates as functions of the discount and explicit conditions
for uniformity and reference continuity as the discount tends to zero.
The proof permits active control constraints and derives optimality
conditions by admissible first variations. The statement is sufficient,
not a necessary-and-sufficient classification of all nonlinear systems.
\end{abstract}

\section{Scope of the question and of this answer}
The canonical C61-93 statement asks for structural assumptions yielding a
discount-dependent reference, an \emph{unweighted} measure-turnpike bound
uniform in the horizon and in a specified data family, and quantitative
dependence on the discount, including sufficient conditions for uniformity
as the discount tends to zero. It expressly does not assert a universal
turnpike under continuity, compactness, Lipschitz dynamics, and attainment
alone. The source \cite{overview}, Section 5, identifies continuous-time
discounting as a research direction; Section 2 explains the horizon and
initial-condition uniformity in the definition.

The theorem below supplies one general, directly checkable sufficient
structural regime and proves \emph{all} these requested conclusions in that
regime. The drift and costs can be nonlinear, the dimensions are arbitrary,
and the control constraint may be active. No reference, solution of a
canonical boundary-value problem, dissipativity inequality along unknown
optimizers, or turnpike estimate is assumed. The result is not a
necessary-and-sufficient classification of every nonlinear optimal-control
problem, and no claim of novelty is made.

\paragraph{Endpoint and terminal-cost family.}
Initial states range over a fixed compact set $X_0$ strictly inside the
state ball $X$ below. The terminal state is free, subject only to the
state constraint. Terminal costs range over a family $\mathcal F$ with
uniform bounds on their values \emph{and their first derivatives}. Thus
the identified structural hypothesis is stronger than merely bounded
terminal values. Prescribed terminal-state constraints are not part of
the displayed C61-93 problem or of the family treated here.

\section{Data assumptions and explicit constants}
Let the state dimension be $d\geq1$ and the control dimension be $r\geq1$.
Use Euclidean vector norms and their induced matrix norms. Let
\[
 X=\{x\in\R^d:|x|\leq R_X\},\qquad R_X>0,
\]
and let $U\subset\R^r$ be nonempty, compact, and convex. Fix a nonempty
compact set $X_0\subset\operatorname{int}X$. Consider
\begin{equation}\label{eq:problem}
 \min_{u(t)\in U}\left\{
 \int_0^T e^{-\rho t}\bigl(L(x(t))+Q(u(t))\bigr)\,dt
       +e^{-\rho T}\Phi(x(T))\right\},\qquad
 \dot x=g(x)+Bu,\quad x(0)=x_0,
\end{equation}
with $x(t)\in X$, arbitrary $T>0$, $x_0\in X_0$, and $\Phi\in\mathcal F$.
Here $\rho\geq0$; the positive values are the discounted problems and
$\rho=0$ is included to make the limiting assertions precise.

Assume the following inequalities hold with finite constants
$a,m,\eta>0$ and $K,M,G,H,F\geq0$:
\begin{enumerate}
\item $g$ is $C^2$ on a neighborhood of $X$, $B$ is a constant real
matrix, and, for all $x,y\in X$,
\begin{align}
 \sym Dg(x)&\preceq-aI,\label{eq:drift}\\
 \|Dg(x)-Dg(y)\|&\leq K|x-y|.\label{eq:jaclip}
\end{align}
\item The state boundary is uniformly strictly inward for \emph{every}
control:
\begin{equation}\label{eq:inward}
 x\cdot(g(x)+Bu)\leq-\eta
 \qquad (|x|=R_X,\ u\in U).
\end{equation}
\item $L$ is $C^1$ on a neighborhood of $X$ and
\begin{equation}\label{eq:statecost}
 |\nabla L(x)|\leq G,\qquad
 |\nabla L(x)-\nabla L(y)|\leq M|x-y|
 \quad(x,y\in X).
\end{equation}
Convexity of $L$ is not required.
\item $Q$ is $C^1$ on a neighborhood of $U$ and is $m$-strongly convex
on $U$, in the explicit sense
\begin{equation}\label{eq:strong}
 Q(v)\geq Q(u)+\nabla Q(u)\cdot(v-u)+\frac m2|v-u|^2
 \qquad(u,v\in U).
\end{equation}
\item Each $\Phi\in\mathcal F$ is $C^1$ on a neighborhood of $X$, with
\begin{equation}\label{eq:terminal}
 \sup_{\Phi\in\mathcal F}\sup_{x\in X}|\Phi(x)|\leq F,
 \qquad
 \sup_{\Phi\in\mathcal F}\sup_{x\in X}|\nabla\Phi(x)|\leq H.
\end{equation}
Convexity of the terminal costs is not required.
\end{enumerate}
These are bounds on the problem data, independent of unknown optimizers
and of the horizon. The strict inward condition can, for example, be
checked by bounding the maximum of $x\cdot(g(x)+Bu)$ on the compact
boundary and control set.

Write $b=\|B\|$, $k=b^2/m$, $D=\diam X=2R_X$, and, for a given discount,
define
\begin{equation}\label{eq:smallgain}
 M_\rho=M+\frac{KG}{a+\rho},\qquad
 c_\rho=kM_\rho,\qquad
 \kappa_\rho=\frac{c_\rho}{a(a+\rho)}.
\end{equation}
The quantitative small-gain condition at discount $\rho$ is
\begin{equation}\label{eq:sgcondition}
 \kappa_\rho<1.
\end{equation}
For this discount put
\begin{equation}\label{eq:rate}
 \gamma_*(\rho)=a+\frac\rho2-
                 \sqrt{\frac{\rho^2}{4}+c_\rho}>0.
\end{equation}
Choose any $0<\gamma<\gamma_*(\rho)$ and define
\begin{align}
 \theta_{\rho,\gamma}
   &=\frac{c_\rho}{(a-\gamma)(a+\rho-\gamma)}<1,
       & H_\rho&=H+\frac G{a+\rho},\label{eq:theta}\\
 A_{\rho,\gamma}
   &=\frac{D+\dfrac{kH_\rho}{a-\gamma}}
           {1-\theta_{\rho,\gamma}},
       & P_{\rho,\gamma}
   &=H_\rho+\frac{M_\rho A_{\rho,\gamma}}{a+\rho-\gamma},\label{eq:AP}\\
 C_{\rho,\gamma}&=A_{\rho,\gamma}+\frac bmP_{\rho,\gamma}.
 \label{eq:C}
\end{align}
Here $\gamma<a$ follows from \eqref{eq:rate}; all denominators are positive.

\begin{theorem}[Global discounted turnpike under data small gain]\label{thm:main}
Under the preceding data assumptions, every problem \eqref{eq:problem}
has an optimal control. For every discount satisfying
\eqref{eq:sgcondition}, a unique stationary current-value reference
$(\bar x_\rho,\bar u_\rho,\bar p_\rho)$ with Hamiltonian-minimizing control
is constructed from the data by the contraction specified below.
For \emph{every} optimal pair and every $T>0$, $x_0\in X_0$,
$\Phi\in\mathcal F$, one has
\begin{equation}\label{eq:exp}
 |x_T(t)-\bar x_\rho|+|u_T(t)-\bar u_\rho|
 \leq C_{\rho,\gamma}
             \left(e^{-\gamma t}+e^{-\gamma(T-t)}\right)
 \quad\text{for almost every }t\in[0,T].
\end{equation}
The optimal control has a continuous representative for which the
inequality holds everywhere. In particular, for every $\varepsilon>0$,
\begin{equation}\label{eq:measure}
 \Leb\{t\in[0,T]:|x_T(t)-\bar x_\rho|
                   +|u_T(t)-\bar u_\rho|>\varepsilon\}
 \leq\frac2\gamma\logplus\left(\frac{2C_{\rho,\gamma}}\varepsilon\right),
\end{equation}
where $\logplus s=\max(0,\log s)$. The constants are independent of $T$,
$x_0$, and the chosen terminal cost in the specified family.

If $c_0<a^2$, choose any fixed
\begin{equation}\label{eq:uniformgamma}
 0<\gamma<a-\sqrt{c_0}.
\end{equation}
Then all discounts $\rho\geq0$ satisfy the theorem with this same rate
and with $C_{\rho,\gamma}\leq C_{0,\gamma}$. Thus both the exponential
and the unweighted measure estimates are uniform as $\rho\downarrow0$.
The reference is Lipschitz in $\rho$ in this uniform regime, with explicit
constants given in Proposition~\ref{prop:continuity} below.
\end{theorem}

\section{Construction of the reference and existence of optimizers}
\begin{lemma}[Invariance, constant-control equilibria, and feedback]\label{lem:maps}
Every measurable $U$-valued control starting in $X_0$ produces a trajectory
strictly inside $X$ for all finite times. For each constant $u\in U$,
there is a unique $z(u)\in\operatorname{int}X$ satisfying
$g(z(u))+Bu=0$, and
\begin{equation}\label{eq:zlip}
 |z(u)-z(v)|\leq\frac ba|u-v|.
\end{equation}
The map
\begin{equation}\label{eq:feedback}
 \mathcal K(p)=\argmin_{u\in U}\{Q(u)+p\cdot Bu\}
\end{equation}
is single valued and satisfies
\begin{equation}\label{eq:feedbacklip}
 |\mathcal K(p)-\mathcal K(q)|\leq\frac bm|p-q|.
\end{equation}
\end{lemma}
\begin{proof}
Compactness and continuity in \eqref{eq:inward} provide a boundary shell
on which $x\cdot(g(x)+Bu)\leq-\eta/2$ uniformly over $u\in U$.
The absolutely continuous function $|x(t)|^2$ therefore has strictly
negative derivative almost everywhere while it lies in that shell.
It cannot cross the shell outward. Starting in the fixed compact
$X_0\subset\operatorname{int}X$ gives a positive uniform distance from
the boundary for all such controls. Local existence and uniqueness of
the controlled ODE, followed by compact continuation, give global
existence. This also proves that every convex control variation is
admissible; state-constraint multipliers will not be needed.

For the same constant control, \eqref{eq:drift} and convexity of $X$ give
\[
 \frac d{dt}|x(t)-y(t)|^2\leq-2a|x(t)-y(t)|^2.
\]
The time-one autonomous flow is therefore a strict contraction of the
complete space $X$ into itself and has a unique fixed point $z$.
For every $s\geq0$, the time-$s$ image of $z$ is again a fixed point of
the time-one flow, by autonomy and commutation of the flow maps.
Uniqueness forces that image to equal $z$, so $z$ is an equilibrium.
Strict inwardness places it in the interior. Equilibria for controls
$u,v$ satisfy
\[
 a|z(u)-z(v)|^2
 \leq B(u-v)\cdot(z(u)-z(v)),
\]
which implies \eqref{eq:zlip}.

The minimum in \eqref{eq:feedback} exists by compactness and is unique
by strong convexity. Its first-order variational inequality is
\[
 [\nabla Q(\mathcal K(p))+B^Tp]\cdot[w-\mathcal K(p)]\geq0
 \qquad(w\in U).
\]
Use $w=\mathcal K(q)$ in this inequality and the reverse choice in
the corresponding inequality for $q$. Strong monotonicity of $\nabla Q$
then yields
$m|\mathcal K(p)-\mathcal K(q)|^2\leq
 b|p-q||\mathcal K(p)-\mathcal K(q)|$, proving the Lipschitz bound.
This argument includes active boundary controls.
\end{proof}

For $x\in X$ define
\begin{equation}\label{eq:stationaryp}
 \mathcal P_\rho(x)=(\rho I-Dg(x)^T)^{-1}\nabla L(x).
\end{equation}
The inverse exists: the symmetric part of the matrix in parentheses is
at least $(a+\rho)I$. In particular,
\begin{equation}\label{eq:plip}
 |\mathcal P_\rho(x)|\leq\frac G{a+\rho},\qquad
 |\mathcal P_\rho(x)-\mathcal P_\rho(y)|
       \leq\frac{M_\rho}{a+\rho}|x-y|.
\end{equation}
For the second inequality, multiply the difference by
$\rho I-Dg(x)^T$. The result is
\[
 \nabla L(x)-\nabla L(y)
       +(Dg(x)-Dg(y))^T\mathcal P_\rho(y),
\]
whose norm is at most $M_\rho|x-y|$.

\begin{lemma}[Reference construction]\label{lem:reference}
If $\kappa_\rho<1$, the data-defined map
\begin{equation}\label{eq:contraction}
 \mathcal T_\rho(u)=\mathcal K(\mathcal P_\rho(z(u)))
\end{equation}
is a contraction of $U$ with factor at most $\kappa_\rho$.
Its unique fixed point defines
\[
 \bar u_\rho=\mathcal T_\rho(\bar u_\rho),\qquad
 \bar x_\rho=z(\bar u_\rho),\qquad
 \bar p_\rho=\mathcal P_\rho(\bar x_\rho).
\]
They satisfy
\begin{align}
 g(\bar x_\rho)+B\bar u_\rho&=0,\label{eq:ref1}\\
 [\nabla Q(\bar u_\rho)+B^T\bar p_\rho]
                 \cdot(u-\bar u_\rho)&\geq0\quad(u\in U),\label{eq:ref2}\\
 \nabla L(\bar x_\rho)+Dg(\bar x_\rho)^T\bar p_\rho
                 &=\rho\bar p_\rho.\label{eq:ref3}
\end{align}
For an interior reference control, \eqref{eq:ref2} is the equality
$\nabla Q(\bar u_\rho)+B^T\bar p_\rho=0$ requested in C61-93.
\end{lemma}
\begin{proof}
Combine \eqref{eq:zlip}, \eqref{eq:feedbacklip}, and \eqref{eq:plip}.
The product of the three Lipschitz constants is exactly
$b^2M_\rho/[ma(a+\rho)]=\kappa_\rho<1$.
Banach's fixed-point theorem gives the fixed point as the limit of
successive iteration of this known map. The displayed equations follow
from the definitions. Conversely, any stationary triple whose control
minimizes the Hamiltonian satisfies this fixed-point equation, so it
coincides with the constructed triple.
\end{proof}

\begin{lemma}[Attainment for the full specified family]\label{lem:existence}
For each finite $T>0$, $x_0\in X_0$, $\rho\geq0$, and
$\Phi\in\mathcal F$, problem \eqref{eq:problem} has an optimizer.
\end{lemma}
\begin{proof}
All measurable $U$-valued controls are feasible by Lemma~\ref{lem:maps}.
Their states stay in a common compact set and have uniformly bounded
derivatives. Take a minimizing sequence. Its controls have a weakly
convergent subsequence in $L^2(0,T;\R^r)$, and its states have a uniformly
convergent subsequence by equicontinuity. The weak limit is $U$-valued
almost everywhere because this is a closed convex subset of $L^2$.
Passing to the integrated state equation uses uniform convergence of
$g(x)$ and weak convergence of $Bu$ against indicator functions of time
intervals. Thus the limit is a feasible state-control pair.
The state running cost and terminal cost converge. The control integral
is weakly lower semicontinuous, since its integrand is continuous and
convex in $u$ and the weight is positive. This last fact also follows
by taking strongly convergent convex combinations of a weakly convergent
sequence and using convexity of $Q$. The limit therefore attains the
infimum. Compactness makes all costs finite, and no convexity of $L$ or
$\Phi$ was needed.
\end{proof}

\section{First variation and the exponential estimate}
\begin{lemma}[A derived normal current-value adjoint]\label{lem:adjoint}
For every optimal pair, let $p$ solve the linear backward equation
\begin{equation}\label{eq:adjoint}
 \dot p=(\rho I-Dg(x)^T)p-\nabla L(x),\qquad
 p(T)=\nabla\Phi(x(T)).
\end{equation}
Then
\begin{equation}\label{eq:optimalfeedback}
 u(t)=\mathcal K(p(t))\quad\text{for almost every }t.
\end{equation}
In particular, this adjoint and its terminal condition are consequences
of ordinary first variation; normality and a uniform terminal-gradient
bound are not being inferred from bounded terminal values.
\end{lemma}
\begin{proof}
For any measurable $U$-valued $v$, the controls
$u+\varepsilon(v-u)$ are admissible for $0\leq\varepsilon\leq1$,
including the state constraint, by Lemma~\ref{lem:maps}.
Their state derivative at $\varepsilon=0$ solves
\[
 \dot z=Dg(x)z+B(v-u),\qquad z(0)=0.
\]
Differentiate the finite-horizon objective and integrate by parts with
the explicitly constructed solution of \eqref{eq:adjoint}. The result is
\[
 \int_0^T e^{-\rho t}
      [\nabla Q(u(t))+B^Tp(t)]\cdot[v(t)-u(t)]\,dt\geq0.
\]
Localized variations, first for a countable dense subset of $U$ and then
by continuity, imply the corresponding pointwise variational inequality
for every $v\in U$, almost everywhere. Strong convexity makes it
equivalent to \eqref{eq:optimalfeedback}. Thus there is no abnormal
multiplier issue: the running-cost coefficient was one in a direct
derivation from optimality. Strict state invariance justifies the
variations without hidden state-boundary measures. Since $p$ is
continuous and $\mathcal K$ is Lipschitz, the feedback provides a
continuous representative of the optimal control.
\end{proof}

\begin{proof}[Proof of Theorem~\ref{thm:main}]
Existence and the reference were proved above. Write
\[
 y=x-\bar x_\rho,\quad v=u-\bar u_\rho,\quad q=p-\bar p_\rho.
\]
Strong contraction of the drift, the state equation, and
\eqref{eq:feedbacklip} give the scalar comparison estimate
\begin{equation}\label{eq:forward}
 |y(t)|\leq e^{-at}|y(0)|
       +k\int_0^t e^{-a(t-s)}|q(s)|\,ds.
\end{equation}
Indeed, the upper derivative of $|y|$ is at most $-a|y|+b|v|$,
and $|v|\leq(b/m)|q|$.

Subtract \eqref{eq:ref3} from \eqref{eq:adjoint}. The resulting equation is
\[
 \dot q=(\rho I-Dg(x)^T)q-d(t),\qquad
 d(t)=\nabla L(x)-\nabla L(\bar x_\rho)
            +(Dg(x)-Dg(\bar x_\rho))^T\bar p_\rho.
\]
By \eqref{eq:plip}, $|d(t)|\leq M_\rho|y(t)|$.
The backward evolution generated by $\rho I-Dg(x(t))^T$ has norm at most
$e^{-(a+\rho)(s-t)}$ from $s$ back to $t\leq s$; reverse time and apply
\eqref{eq:drift} to see this. Also
$|q(T)|\leq H+G/(a+\rho)=H_\rho$. Hence
\begin{equation}\label{eq:backward}
 |q(t)|\leq H_\rho e^{-(a+\rho)(T-t)}
       +M_\rho\int_t^T e^{-(a+\rho)(s-t)}|y(s)|\,ds.
\end{equation}
These estimates have been derived for every optimizer, not assumed for a
solution of a separate boundary-value problem.

Set $E_T(t)=e^{-\gamma t}+e^{-\gamma(T-t)}$, and define the finite norms
\[
 Y=\max_{0\leq t\leq T}\frac{|y(t)|}{E_T(t)},\qquad
 P=\max_{0\leq t\leq T}\frac{|q(t)|}{E_T(t)}.
\]
Elementary integration gives
\begin{align*}
 \int_0^t e^{-a(t-s)}E_T(s)\,ds
     &\leq \frac{E_T(t)}{a-\gamma},\\
 \int_t^T e^{-(a+\rho)(s-t)}E_T(s)\,ds
     &\leq \frac{E_T(t)}{a+\rho-\gamma}.
\end{align*}
For example, the two forward summands have the respective bounds
$e^{-\gamma t}/(a-\gamma)$ and
$e^{-\gamma(T-t)}/(a+\gamma)$; the backward ones are analogous.
Since $|y(0)|\leq D$, equations \eqref{eq:forward}--\eqref{eq:backward}
therefore imply
\[
 Y\leq D+\frac{k}{a-\gamma}P,
 \qquad
 P\leq H_\rho+\frac{M_\rho}{a+\rho-\gamma}Y.
\]
The product of the coefficients is $\theta_{\rho,\gamma}<1$.
Solving these two inequalities gives $Y\leq A_{\rho,\gamma}$ and
$P\leq P_{\rho,\gamma}$. Finally
$|v|\leq(b/m)|q|$ proves \eqref{eq:exp} with exactly the constant
\eqref{eq:C}.

The admissible range of $\gamma$ follows by solving
$(a-\gamma)(a+\rho-\gamma)>c_\rho$ with $0<\gamma<a$;
its smaller root is \eqref{eq:rate}. To obtain the measure estimate,
let $s=\gamma^{-1}\logplus(2C_{\rho,\gamma}/\varepsilon)$.
Outside the two endpoint intervals of length $s$, both exponentials in
\eqref{eq:exp} are at most $\varepsilon/(2C_{\rho,\gamma})$.
The exceptional set has ordinary Lebesgue measure at most $2s$,
including when the intervals overlap. This argument uses no discounted
measure.

If $c_0<a^2$ and \eqref{eq:uniformgamma} holds, then $M_\rho,c_\rho,H_\rho$
are nonincreasing in $\rho$, whereas $a+\rho-\gamma$ increases.
Consequently
\[
 \theta_{\rho,\gamma}\leq\frac{c_0}{(a-\gamma)^2}<1,
 \quad A_{\rho,\gamma}\leq A_{0,\gamma},
 \quad P_{\rho,\gamma}\leq P_{0,\gamma}.
\]
This proves the uniform statements; reference continuity is proved next.
\end{proof}

\section{Discount dependence of the reference}
\begin{proposition}[Explicit uniform reference continuity]\label{prop:continuity}
Assume $c_0<a^2$, write $\kappa_0=c_0/a^2$, and put
\[
 L_u=\frac{bG}{m a^2(1-\kappa_0)},\qquad
 L_x=\frac baL_u,\qquad
 L_p=\frac{M_0}{a}L_x+\frac G{a^2}.
\]
For all $\rho,\sigma\geq0$,
\[
 |\bar u_\rho-\bar u_\sigma|\leq L_u|\rho-\sigma|,
 \quad |\bar x_\rho-\bar x_\sigma|\leq L_x|\rho-\sigma|,
 \quad |\bar p_\rho-\bar p_\sigma|\leq L_p|\rho-\sigma|.
\]
In particular the reference converges to the zero-discount reference
without any horizon-dependent or terminal-cost-dependent choice.
\end{proposition}
\begin{proof}
The resolvent identity at a fixed $x$ gives
\[
 |\mathcal P_\rho(x)-\mathcal P_\sigma(x)|
 \leq\frac{G|\rho-\sigma|}{(a+\rho)(a+\sigma)}.
\]
Compare the fixed points of \eqref{eq:contraction}, using one map's
contraction bound and then this estimate for the difference between the
two maps at a fixed control. Since $\kappa_\rho\leq\kappa_0<1$, this yields
the asserted $L_u$. Equation \eqref{eq:zlip} gives $L_x$, and
\eqref{eq:plip} together with the resolvent identity gives $L_p$.
\end{proof}

\paragraph{Interpretation of the rate and discount restriction.}
The bound \eqref{eq:rate} is a guaranteed admissible rate, not a claim of
sharpness or a necessary condition. All of its discount dependence is
explicit. Equivalently, writing $k=b^2/m$ and $y=a+\rho$, the small-gain
condition is
\[
 a y^2-kMy-kKG>0.
\]
This directly specifies a sufficient discount range from the data.
For vanishing-discount uniformity the strict margin $c_0<a^2$ suffices.
Failure of that sufficient inequality does not establish absence of a
turnpike.

\paragraph{Why the reference generally moves with the discount.}
For illustration within the theorem, consider scalar
$g(x)=-ax$, $B=b$, $L(x)=(x-z)^2/2$, and $Q(u)=mu^2/2$, with a sufficiently
large compact control interval and invariant state ball, and $m>b^2/a^2$.
The reference is
\[
 \bar u_\rho=\frac{abz}{ma(a+\rho)+b^2},\qquad
 \bar x_\rho=\frac{b^2z}{ma(a+\rho)+b^2}.
\]
It differs from the undiscounted static minimizer whenever
$\rho>0$ and $bz\ne0$. This example illustrates the reference formula;
the theorem and its proof cover the nonlinear, nonquadratic class above.

\section{Why the base assumptions alone are insufficient}
Take $X=\{x\in\R^2:|x|\leq1\}$, $U=\{0\}$, $L=Q=\Phi=0$, and
\[
 \dot x=Jx,\qquad
 J=\begin{pmatrix}0&-1\\1&0\end{pmatrix},\qquad x(0)=(1/2,0).
\]
All data are smooth, the constraints are compact, the dynamics are
uniformly Lipschitz, and for every horizon and discount the unique
admissible state-control trajectory is optimal. Terminal costs are
uniformly zero. The only stationary pair is $(\bar x,\bar u)=(0,0)$.
The unique trajectory has $|x(t)|=1/2$ for every $t$, so for
$\varepsilon=1/4$ its unweighted exceptional set has measure $T$.
No horizon-independent bound exists. This is a full counterexample to
the \emph{base-assumptions-only} implication, which C61-93 itself does
not assert. It identifies why additional structural assumptions are
being supplied rather than silently presumed.

\section{Audit checklist and limits of the claim}
\begin{itemize}
\item The reference is constructed by a data-defined contraction and
obeys the discounted current-value stationary equations (with the correct
variational inequality for an active control constraint).
\item Optimizer existence and the adjoint are proved; the argument does
not start by assuming a suitable canonical-system solution.
\item Strict state invariance is required for every control and removes
the state-constraint issue in the first-variation derivation.
\item The terminal family is uniformly bounded in $C^1$; bounded values
alone do not imply the gradient estimate used in the proof.
\item The conclusion is for every optimizer, every positive horizon,
all stated initial data and terminal costs, and ordinary Lebesgue measure.
\item Both discount-dependent constants and explicit uniformity
conditions as $\rho\downarrow0$ are given.
\item The theorem is a sufficient structural answer to the stated
agenda. It is not a classification of every nonlinear drift, a theorem
for hard terminal-state constraints, or a claim of research novelty.
\end{itemize}

\section*{Authorship and verification}
Prepared by Ji Ho Bae with assistance from GPT Codex. The manuscript
received internal symbolic and scope checks. It has not been formally
verified in Lean and has not undergone external peer review. No claim of
novelty is made.

\noindent\textbf{Final scope: Sufficient structural answer to the formal agenda.}
The result supplies the requested sufficient hypotheses, constructed
reference, unweighted bound, discount-dependent rates and constants, and
vanishing-discount uniformity for the stated general nonlinear class.
The qualifications and endpoint-family specification above are part of
the theorem.

\begin{thebibliography}{9}
\bibitem{overview}
T. Faulwasser and L. Gr\"une,
\emph{Turnpike Properties in Optimal Control: An Overview of Discrete-Time
and Continuous-Time Results}, arXiv:2011.13670v3, 2021.
Sections 2 and 5.
\url{https://arxiv.org/html/2011.13670v3}.
\end{thebibliography}
\end{document}
